\begin{abstract}
We position Goal-Frontier Maximization as a measure-theoretic formalization of
    the capability welfare economics tradition of Sen, Nussbaum, Becker, and
    Stiglitz--Sen--Fitoussi: a tradition rich in qualitative content and in
    counting-based quantitative measurement (Alkire--Foster) but without a
    unified measure-theoretic substrate that admits cooperative excess and
    polarity discipline as first-class structure.  The formalization yields two results of
    structurally distinct kinds: a \emph{fungibility-collapse correspondence
    theorem} establishing GFM as a strict superset of standard utility theory
    at the linear / risk-neutral case, and a \emph{Goodhart theorem} bounding
    each Lipschitz alignment property's deviation from the operational truth
    (the realized exercise contribution $T = \volRwin$) by $\mathrm{Lip}(g)
    \cdot \epsnonres$, with property-specific Lipschitz analyses left to
    subsequent work.  The proxy-truth gap admits four
    operational-intervention channels for tightening (observation density,
    attestation quality, individuation discipline, bundle decomposition),
    each formalizing a move the qualitative tradition has been negotiating
    informally.  Admissible improvements to the economic
    model tighten certified alignment bounds in the direction of the
    operational truth, above the deployment-specific residual-class floor.
\end{abstract}
